\section{The complete finite-memory frontier for reset control}
\label{sec:r78-reset-frontier}

The optimal-memory threshold and the cost at every smaller memory are
different invariants.  We now determine the latter by an attained finite
variational principle.  The result applies to the full reset class of
Section~\ref{sec:reset-control}; strict positivity of the failure effects
is not required.  In particular it includes reports which a least-favorable
preparation cannot reach.  The key reduction preserves the controller's
existing command rows while removing all dependence of the post-report
update on the preceding label and command.

Keep the primitive data $E_{a,j},G_a,\sigma_i$, the positive cost floor,
and the absence of a closed failure set from
Theorem~\ref{thm:reset-control}.  For a command distribution
$p\in\Delta(\mathcal A)$ define
\[
 c_i(p)=\sum_a p_a c_{ia},\qquad
 Q_{ij}(p)=\sum_a p_a Q_{iaj},\qquad
 G(p)=\sum_a p_aG_a,\qquad E_j(p)=\sum_a p_aE_{a,j}.
\]
All these quantities concern the actual reported reset state.  There is
no additional measurement, cost observation, free clock, or access to the
unknown initial density matrix.

Let $V_M$ be the infimum of the worst-initial-state expected cost over
autonomous controllers with at most $M\ge1$ live labels, a prescribed
pure initial label, and mandatory stopping at success.  The command row
at a label can be mixed.  Updates after a failure may initially depend
stochastically on the old label, the sampled command, and its report.
Controllers with infinite worst-state cost do not affect the infimum.

\subsection{A memory-preserving deterministic update theorem}

\begin{theorem}[Dominance with the same command rows]
\label{thm:r78-palette-normal-form}
Suppose an $M$-label controller has finite worst-initial-state cost $C$,
and its command rows are $p^1,\ldots,p^M$, with pure initial label $1$.
There is another controller with these same rows and the same initial
label whose update after report $i$ is one deterministic label $k(i)$,
independent of the preceding label and command.  Its expected cost is
no greater than the original cost at every initial density matrix.

More precisely, the finite system
\begin{equation}\label{eq:r78-palette-bellman}
 x_i=\min_{1\le k\le M}
       \left\{c_i(p^k)+\sum_jQ_{ij}(p^k)x_j\right\}
\end{equation}
has a unique finite nonnegative solution, with $0\le x_i\le C$.
Choose $k(i)$ attaining the minimum.  The resulting continuation is
proper from every reset state, and its full cost from an initial $\rho$
is exactly
\begin{equation}\label{eq:r78-palette-start}
 \tr\bigl(\rho H(p^1,x)\bigr),\qquad
 H(p,x)=G(p)+\sum_jx_jE_j(p).
\end{equation}
This is a pointwise cost comparison, without a convexification of the
hard memory constraint.
\end{theorem}
\begin{proof}
Regard the finite list of rows as the available choices at each reset
state.  A choice $k$ has mean cost $c_i(p^k)$ and substochastic reset
transition row $Q_{ij}(p^k)$.  Any extra information about the sampled
command or earlier private coins does not alter the physical state
$\sigma_j$ after report $j$.

The original controller can be started with label $1$ and physical
state $\sigma_i$, for every $i$: the robustness quantifier includes
\emph{all} initial density matrices.  This gives, for the same list
of rows, a proper history-dependent policy of cost at most $C$ from
each reset state.  This comparison restart is used only to prove
finiteness; it is not an additional operation supplied to the eventual
controller.  In particular it remains valid when some reset reports
are inaccessible from a least-favorable preparation.

Starting the finite-horizon recursion for the displayed row choices
at zero gives increasing vectors bounded above by $C\mathbf1$.
Their limit $x$ is finite and satisfies
\eqref{eq:r78-palette-bellman}.  For a minimizing selector $k(i)$,
iterate its equality through a finite horizon.  The accumulated
expected cost is bounded by $x_i$.  The positive cost floor therefore
bounds the sum of survival probabilities, so the selector is proper
from every state.  The bounded remaining potential tends to zero;
its expected total cost is $x_i$.

Conversely Bellman's inequality holds conditionally at each call of
any policy using these rows, even if it retains the sampled command
and earlier coins.  Telescope it to a finite horizon.  For a
finite-cost policy, the survival remainder vanishes by the same
positive floor, proving that its cost from reset state $i$ is at
least $x_i$.  Infinite-cost policies cannot improve the bound.  The
same verification applied to any finite nonnegative solution gives
uniqueness.

After the original first call, each reached continuation from report
$j$ uses the same list of rows and hence costs at least $x_j$.
Conditioning on that call proves that its cost at every $\rho$ is
at least \eqref{eq:r78-palette-start}.  The new controller starts
with row $p^1$ and, after report $j$, writes precisely the existing
label $k(j)$.  Thereafter it implements the minimizing selector and
attains this expression.  The actual device thus uses at most the
same $M$ labels, with no uncounted row-selection memory.
\end{proof}

The all-initial-state quantifier is essential to the bounded comparison
in this proof.  Properness for just one supplied initial state would
not by itself imply a finite continuation value at every reset state.
Nor does the theorem assert that mixed command rows can always be
replaced by pure commands: only their post-report selection is made
deterministic.

\subsection{A compact exact variational principle}

Choose once and for all a full-support row $\bar p$, for example the
uniform row on $\mathcal A$.  The proof of
Theorem~\ref{thm:reset-control} makes $\bar Q=(Q_{ij}(\bar p))$
transient.  Consequently
\begin{equation}\label{eq:r78-explicit-cap}
 \bar x=(I-\bar Q)^{-1}(c_i(\bar p))_{i=1}^s,
 \qquad C_0=\lambda_{\max}H(\bar p,\bar x)<\infty.
\end{equation}
The constant one-row controller has worst-state cost $C_0$, and
$\bar x_i=\tr(\sigma_iH(\bar p,\bar x))\le C_0$.
Thus this is a primitive, explicit compactness bound, not an assumed
upper bound on an unknown optimal trajectory.

For $k:\{1,\ldots,s\}\to\{1,\ldots,M\}$ let $\mathcal F_{M,k}$
consist of the rows $p^1,\ldots,p^M\in\Delta(\mathcal A)$ and
$x\in[0,C_0]^s$ satisfying
\begin{equation}\label{eq:r78-supersolution}
 x_i\ge c_i(p^{k(i)})+\sum_jQ_{ij}(p^{k(i)})x_j
 \qquad(1\le i\le s).
\end{equation}
There are $M^s$ possible selectors and each corresponding feasible
set is compact; some of them may be empty.

\begin{theorem}[Exact attained frontier at every memory]
\label{thm:r78-reset-frontier}
For every $M\ge1$,
\begin{equation}\label{eq:r78-full-frontier}
 V_M=\min_k\ \min_{(p^1,\ldots,p^M,x)\in\mathcal F_{M,k}}
             \lambda_{\max}H(p^1,x).
\end{equation}
The minimum is finite and attained, and an attaining controller has
deterministic report-dependent updates as in
Theorem~\ref{thm:r78-palette-normal-form}.  For $0\le b\le C_0$,
there is an $M$-label controller of worst-state cost at most $b$ if
and only if one of the finite systems
\begin{equation}\label{eq:r78-budget-certificate}
 (p^1,\ldots,p^M,x)\in\mathcal F_{M,k},\qquad
 H(p^1,x)\preccurlyeq bI
\end{equation}
is feasible.

In particular $V_M$ is nonincreasing in $M$ and equals the unrestricted
value $V_*$ once $M\ge\min(s,|\mathcal A|)+1$.  The exact optimal
label count, without a strict-exposure hypothesis, is
\begin{equation}\label{eq:r78-general-threshold}
 m_{\mathrm{opt}}=
 \min\{M:\text{\eqref{eq:r78-budget-certificate} is feasible at }b=V_*\}.
\end{equation}
Every $M<m_{\mathrm{opt}}$ has the strict gap $V_M>V_*$.
For algebraic primitive data, $V_M$, $m_{\mathrm{opt}}$, and an
attaining controller admit finite exact real-algebraic procedures.
No polynomial complexity or general minimax saddle at fixed $M$ is
asserted.
\end{theorem}
\begin{proof}
A feasible supersolution bounds every finite-prefix continuation
cost under the displayed selector by $x_i$.  The positive cost floor
therefore makes that selector proper from every reset state.  Its
true continuation vector is at most $x$, and its full expected cost
operator is at most $H(p^1,x)$, because every $E_j(p^1)$ is positive
semidefinite.  Thus every feasible point is an actual controller
upper certificate.

Conversely take any controller of worst-state cost $b\le C_0$.
Theorem~\ref{thm:r78-palette-normal-form} supplies its row-list
Bellman vector $x\le b\mathbf1$, a minimizing selector, and a
cost operator bounded by $bI$.  These are feasible in
\eqref{eq:r78-budget-certificate}.  The bound $x\le b\mathbf1$
uses the comparison starting at each $\sigma_i$ with label $1$;
it is not inferred from positive exposure of a particular dual state.

All rows equal to $\bar p$, with $x=\bar x$, give a feasible point
of value $C_0$, for every $M$.  A finite union of the compact
$\mathcal F_{M,k}$ is therefore nonempty and the continuous
objective attains its minimum.  If the controller infimum is strictly below $C_0$, apply the converse
to an approximating sequence below $C_0$; if it equals $C_0$, use
the displayed one-row certificate.  Together with the first implication,
this identifies the minimum with the controller infimum.  At a minimizing feasible
point the constructed controller cannot have strictly smaller
worst-state cost, since that would contradict the equality of the
infima.  This proves actual attainment as well as the budget
equivalence.

Allowing duplicate rows shows monotonicity.  The unrestricted
implementation in Theorem~\ref{thm:reset-control} proves the stated
finite saturation bound.  At every smaller $M$, equality with $V_*$
is equivalent to an attaining optimal controller by the result
already proved, giving \eqref{eq:r78-general-threshold} and the
strict-gap assertion.  This supplies a general variational
classification even where the simpler transversal formula of
Theorem~\ref{thm:r77-reset-cover} is inapplicable.

For algebraic data, $\bar x$ and $C_0$ are algebraic.  For each
selector, the row constraints and
\eqref{eq:r78-supersolution} are polynomial inequalities with
algebraic coefficients.  On separating real and imaginary parts,
the Hermitian inequality in \eqref{eq:r78-budget-certificate} is
equivalent to nonnegativity of its principal minors.  Exact
real-algebraic elimination determines its feasible budget set and
least endpoint; compactness ensures that endpoint is included and
provides algebraic attaining coordinates.  Finally enumerate the
finite range in \eqref{eq:r78-general-threshold}.  This is an exact
computability conclusion from the proved finite system, not a
claim that a numerical optimization certifies a tie or that an
elimination computation has been performed for arbitrary inputs.
\end{proof}

With the inherited single terminal decoder, event stopping has the
same total-label frontier $V_M$.  A state-stopped device with $M+1$
total labels has the live-label frontier $V_M$, since its terminal
label is disjoint.  The prescribed pure start is already charged
in \eqref{eq:r78-full-frontier}: $p^1$ is the command distribution
at a counted initial label, not a freely mixed choice of hidden
initial labels.

\begin{proposition}[Continuity of the full frontier]
\label{prop:r78-frontier-stability}
Fix the Hilbert dimension and the finite command and report alphabets.
For every fixed $M$, $V_M$ is continuous in the primitive effects,
costs, and reset states on the set of valid reset instruments with a
positive cost floor and no closed failure set.  Thus if $V_M>V_*$ at
one instrument, the same strict inequality holds throughout some
neighborhood within this class.  This assertion concerns values and
strict gaps; the exact optimal-memory count need not be continuous.
\end{proposition}
\begin{proof}
Near a fixed instrument, the same full-support row $\bar p$ has
transient $\bar Q$: its spectral radius is strictly less than one,
an open condition.  The resolvent in
\eqref{eq:r78-explicit-cap} and its spectral value $C_0$ are
continuous, so the compactness bounds are locally uniform.

An attaining controller can be chosen with deterministic report
updates and with its true continuation vector finite from every
reset state.  Its substochastic reset matrix $Q^k$, whose $i$th
row is $Q_{ij}(p^{k(i)})$, is transient.  Indeed its finite expected
cost and the positive cost floor bound expected duration from every
state, so the entries of $\sum_{n\ge0}(Q^k)^n$ are finite.
Keep its rows and selector fixed under a small perturbation.  The
same matrix remains transient, and both
$(I-Q^k)^{-1}(c_i(p^{k(i)}))_i$ and the resulting initial cost
operator vary continuously.  This proves upper semicontinuity of
$V_M$.

For the reverse inequality, take a convergent sequence of primitive
data and attaining points in \eqref{eq:r78-full-frontier}, along
a subsequence realizing the lower limit of the values.  The locally
uniform cap bounds the continuation vectors, the rows lie in compact
simplexes, and a further subsequence has a constant selector.  Pass
to a limit in the supersolution inequalities and in the initial
matrix upper bound at their attaining values.  The limit is an
actual controller certificate by the first part of the preceding
proof.  Its value is no greater than the lower limit, proving lower
semicontinuity.  This argument does not need the limiting continuation
vector to lie below one particular choice of cap; its finite
supersolution already constructs a proper controller.

The unrestricted value equals $V_{\min(s,|\mathcal A|)+1}$, so it
is continuous as well.  Continuity of their difference proves the
strict-gap statement.
\end{proof}

The domination step is specific to observed resetting: after each
report the physical state no longer depends on the earlier record.
Bellman verification, finite-state-controller optimization, spectral
minimax at unrestricted memory, and real-algebraic elimination are
classical ingredients
\cite{BertsekasTsitsiklis91,Meuleau99,PoupartBoutilier03,AmatoBernsteinZilberstein10,Seidenberg54}.
The present result combines them in an exact
counted-memory normal form and compact necessary-and-sufficient
budget certificate.  It does not prove such a reduction for a
general partially observed process or for the original squared-sign
risk frontier.

\subsection{Closed frontiers for arbitrary action-cover instances}

The preceding variational formula can be solved explicitly for the
rational family in Theorem~\ref{thm:r77-reset-universality}.  Let
$A_1,\ldots,A_s$ be nonempty proper subsets of $[q]=\{1,\ldots,q\}$,
$q\ge2$.  On the basis $e_1,\ldots,e_s,f_1,\ldots,f_q$, keep
\[
 \sigma_i=|e_i\rangle\langle e_i|,\quad E_{a,i}=\frac rs I,
 \quad G_a(e_i)=1+\mathbf1_{\{a\notin A_i\}},\quad
 G_a(f_j)=3+\mathbf1_{\{a=j\}},\qquad 0\le r<1.
\]
Here $G_a(e_i)$ and $G_a(f_j)$ denote the corresponding diagonal
entries.  The success effect is $(1-r)I$.  This also defines the
degenerate endpoint $r=0$, where every first call succeeds.  For
$H\subseteq[q]$ put
\[
 U(H)=\{i:H\cap A_i=\varnothing\},\qquad
 b_a(H)=|\{i\in U(H):a\notin A_i\}|,
 \qquad \lambda=\frac{r}{s(1-r)}.
\]
Let $b_{(1)}(H)\le\cdots\le b_{(q)}(H)$ be these $q$ nonnegative
integers in increasing order, with repetitions.

\begin{theorem}[The complete action-cover cost frontier]
\label{thm:r78-cover-frontier}
For every $M\ge1$, the exact worst-state cost at $M$ live labels is
\begin{equation}\label{eq:r78-cover-spectrum}
 V_M=3+\frac r{1-r}+
 \min_{\substack{H\subseteq[q]\\ |H|\le M-1}}
 \ \min_{1\le k\le q}
       \frac{1+\lambda\sum_{j=1}^k b_{(j)}(H)}{k}.
\end{equation}
An optimum uses pure continuing command rows indexed by $H$ and an
initial row uniform on $k$ commands with the $k$ smallest values of
$b_a(H)$.  The initial row and these continuing rows suffice for
deterministic report-dependent routing.  Both minima are finite
enumerations, including all tied and degenerate cases.

For a fixed set system, after subtracting $3+r/(1-r)$ the frontier
is a continuous, nondecreasing, piecewise affine concave function
of $\lambda\ge0$.  At $r=0$ all memory values equal $3+1/q$ and
one label is optimal; for $r>0$ the unrestricted-memory threshold
is the transversal number plus one, as in
Theorem~\ref{thm:r77-reset-universality}.
\end{theorem}
\begin{proof}
Fix any list of rows and put
\[
 h_i=\min_{1\le k\le M}p^k(A_i^c),\qquad
 \overline h=\frac1s\sum_i h_i.
\]
The reset probabilities are independent of the command, so
\eqref{eq:r78-palette-bellman} has the explicit solution
\[
 x_i=1+h_i+\frac{r(1+\overline h)}{1-r},
 \qquad \frac1s\sum_i x_i=\frac{1+\overline h}{1-r}.
\]
Its first-cost eigenvalues on $f_j$ are
$3+p^1_j+r(1+\overline h)/(1-r)$; those on $e_i$ are at most
$2+r(1+\overline h)/(1-r)$.  Hence the former strictly dominate,
and the exact cost for these rows after optimal routing is
\begin{equation}\label{eq:r78-cover-row-objective}
 3+\frac r{1-r}+\max_a p^1_a+
               \lambda\sum_i\min_k p^k(A_i^c).
\end{equation}
This formula applies to all row choices; success has probability
$1-r>0$ on every call, so no properness exception is present.

For each row $p^k$ with $k\ge2$, with the other rows fixed, the
last sum in \eqref{eq:r78-cover-row-objective} is a continuous
concave function of $p^k$: each summand is the minimum of one
linear function and constants.  Its minimum over a simplex has a
vertex minimizer.  Start from a global minimum, which exists by
compactness, and replace these rows successively by such vertices.
No replacement can increase the objective; global minimality then
keeps it unchanged.  Thus all continuing rows may be pure.

Let $H$ be their distinct actions.  When $i\notin U(H)$ some
continuing row has zero failure penalty and $h_i=0$.  When
$i\in U(H)$ all continuing rows have penalty one, giving
$h_i=p^1(A_i^c)$, also when $H$ is empty.  It follows that
\[
 \sum_i h_i=\sum_a b_a(H)p^1_a.
\]
The remaining initial-row minimization is therefore
\begin{equation}\label{eq:r78-initial-cap-problem}
 \min_{p\in\Delta([q])}\left\{\max_a p_a+
                                  \lambda\sum_a b_a(H)p_a\right\}.
\end{equation}

For completeness introduce $t\in[1/q,1]$ and impose $p_a\le t$.
For fixed $t$, the linear term is minimized by filling the smallest
$b_a(H)$ coordinates first, each up to $t$, with at most one
partially filled coordinate.  An exchange of mass from a larger
coefficient to a smaller unsaturated one proves this assertion,
including ties.  On each interval
$[1/(k+1),1/k]$ the minimized expression
$t+\lambda\sum_a b_a(H)p_a$ is affine in $t$.  Its minimum has an
endpoint minimizer; the endpoints are exactly $t=1/k$,
$1\le k\le q$.  At such an endpoint a minimizing $p$ is uniform
on $k$ smallest coefficients, and its value is
$(1+\lambda\sum_{j=1}^k b_{(j)}(H))/k$.  The epigraph problem in
$t,p$ has the same optimum as \eqref{eq:r78-initial-cap-problem},
since one may always reduce $t$ to $\max p_a$.  This proves
\eqref{eq:r78-cover-spectrum} and the attaining construction.

For fixed $H,k$ the displayed value after subtracting the common
term is affine with nonnegative slope in $\lambda$.  The finite
minimum is therefore continuous, piecewise affine, concave and
nondecreasing.  At $r=0$ its minimum is $1/q$, attained by a
single uniform row.  For $r>0$, an optimal initial row at the
unrestricted value must be uniform on all $q$ actions; because
each $A_i$ is proper it has positive penalty at every uncovered
state.  Thus zero excess cost requires $H$ to meet every $A_i$,
and a smallest transversal is sufficient.  Equivalently this is
the already proved positive-effect threshold.
\end{proof}

\begin{corollary}[Arbitrarily long strict memory hierarchies]
\label{cor:r78-strict-hierarchy}
For every $n\ge2$ there is a rational reset experiment in dimension
$2n$, with $n$ commands, $n$ failure reports, success probability
$1/2$ at every call and all costs in $[1,4]$, for which
\begin{equation}\label{eq:r78-strict-hierarchy}
 V_M=4+\frac1n+
       \frac{(n-1)\max\{n-M+1,0\}}{n^2}
       \qquad(M\ge1).
\end{equation}
In particular $V_1>V_2>\cdots>V_{n+1}=V_*$, and each of the
first $n$ memory increments saves exactly $(n-1)/n^2$ in worst-state
expected cost.  All failure effects are positive definite.  In this
family the retained first-residual certificate satisfies
\begin{equation}\label{eq:r78-certificate-gap}
 D_M=\frac{V_M-V_*}{2}.
\end{equation}
Thus it is a strict lower certificate below the optimum threshold,
whereas \eqref{eq:r78-strict-hierarchy} is the full repeated-experiment
cost frontier.

The $n$ strict inequalities persist in an open neighborhood within
the class of valid instruments.  They also persist under explicit
arbitrarily small perturbations for which a failure effect and a
cost operator do not commute.  These perturbations retain positive
failure effects, dimension $2n$, and mean cost operators
$I\preccurlyeq G_a\preccurlyeq4I$; the exact
formula \eqref{eq:r78-strict-hierarchy} refers to the unperturbed
instrument.
\end{corollary}
\begin{proof}
Use the preceding construction with $q=s=n$, $A_i=\{i\}$ and
$r=1/2$, so $\lambda=1/n$.  For a continuing action set $H$,
write $u=n-|H|$ for the number of uncovered singleton sets.
Then $b_a(H)=u-1$ on the $u$ uncovered actions and $b_a(H)=u$
on the remaining actions.  If $u>0$, the candidate initial-row
value at a support size $k$ is
\[
 \begin{cases}
  k^{-1}+(u-1)/n,&1\le k\le u,\\[2pt]
  u/n+(1-u/n)/k,&u\le k\le n.
 \end{cases}
\]
The first expression is minimized at $k=u$; the second is
nonincreasing in $k$ and agrees with the first at $u$.  Consequently
a minimum always occurs at $k=n$, with value
$1/n+u(n-1)/n^2$.  For $u=0$ the same value is $1/n$.
Taking the largest allowable $H$ makes
$u=\max\{n-M+1,0\}$ and proves
\eqref{eq:r78-strict-hierarchy}.  The construction uses a uniform
initial command row and one pure row for each action of $H$;
reports in $H$ select their pure row and the remaining reports
select the initial row.  The stated data and positivity properties
follow directly from the displayed primitive construction.

Finally, the Bellman residuals for this family are
$d_{ia}=\mathbf1_{\{a\ne i\}}$.  The matrix in
\eqref{eq:r77-reset-gap} retains only the first post-failure
penalty, so its minimized eigenvalue is obtained from
\eqref{eq:r78-cover-row-objective} by replacing $\lambda=1/n$
by $r/s=1/(2n)$.  The same row-purification argument applies.
For $k\le u$ the candidate is $1/k+(u-1)/(2n)$; for $k\ge u$
it is $u/(2n)+(1-u/(2n))/k$.  Again $k=n$ is a minimizer,
and subtracting $V_*=4+1/n$ gives
$D_M=u(n-1)/(2n^2)$, proving
\eqref{eq:r78-certificate-gap}.  This equality is specific to
the present family and is not a universal identity for $D_M$.

Continuity from Proposition~\ref{prop:r78-frontier-stability},
applied to the finitely many adjacent gaps, proves the neighborhood
assertion.  For an explicit perturbation put
$X=|e_1\rangle\langle f_1|+|f_1\rangle\langle e_1|$ and replace
only $E_{1,1}$ by $E_{1,1}+\epsilon X$, where
\[
 0<\epsilon<\frac{n-1}{8n^2}.
\]
The failure effect remains positive definite, and the corresponding
success effect $\tfrac12I-\epsilon X$ remains positive definite.
All other primitive data are unchanged.  Since
$G_1(e_1)=1$ and $G_1(f_1)=4$, this modified failure effect does
not commute with $G_1$.

These cost and report effects have a joint physical implementation.
Charge one on every failure.  Split the success effect $B_a$ into
\[
 F_{a,9}=(G_a-I)/8,\qquad F_{a,1}=B_a-F_{a,9},
\]
charging nine and one, respectively, and retaining only the common
success report at the control port.  Both effects are positive:
$0\preccurlyeq F_{a,9}\preccurlyeq3I/8$ and
$B_a\succeq(1/2-\epsilon)I$, with $\epsilon<1/32$.
Their sum is $B_a$, and the aggregate mean cost is
$I+8F_{a,9}=G_a$.  Each failure branch uses its stated
measure-and-prepare reset map, and each success subbranch can use
any terminal prepared state.  Thus the perturbation is an actual
instrument with positive, bounded paid costs and no additional
observation, rather than an unsupported independent assignment
of noncommuting cost and report operators.

The perturbation has zero expectation in every $\sigma_i$, so it
changes none of the reset transition probabilities or continuation
Bellman vectors for any chosen list of command rows.  Each such
vector satisfies $x_i\le4$, since reset costs are at most two and
each subsequent call succeeds with probability $1/2$.  Its initial
cost operator changes only by $\epsilon p^1_1x_1X$, of norm at
most $4\epsilon$.  Hence the optimized values satisfy
$|V_M(\epsilon)-V_M(0)|\le4\epsilon$ for every $M$.
Each adjacent gap is therefore at least
$(n-1)/n^2-8\epsilon>0$.  The common unrestricted upper memory
bound still gives $V_{n+1}(\epsilon)=V_*(\epsilon)$.  This proves
the claimed noncommuting realization without extending the
unperturbed closed formula to perturbed data.
\end{proof}

The family shows that the full frontier carries information absent
from its saturation threshold: each intermediate memory can have a
different, attained optimal value under a fixed success probability
and uniformly bounded primitive costs.  Formula
\eqref{eq:r78-cover-spectrum} retains the combinatorial difficulty
of the arbitrary incidence table, while giving a complete finite
answer rather than an efficient algorithm claim.  The reset
assumption and the declared observation port remain essential;
none of these statements completes the separate general stopping,
phase-type timing, or squared-sign minimax-risk problems.
